\documentclass[leqno]{jsarticle}
\usepackage{amssymb}
\usepackage{amsthm}
\usepackage{amsmath}
\usepackage{comment}
	%可換図式用
		%\usepackage[all]{xy}
	%重くなるので使わいときはコメント化しておく。
%定理環境 thmはあまり使わずpropをなるべく使う。
%主張は基本的に証明の中で使い、そのため番号を付けない。
%注意環境を付けてみた。
\newtheorem{thm}{定理}
\newtheorem{prop}[thm]{命題}
\newtheorem{cor}[thm]{系}
\newtheorem{lem}[thm]{補題}
\newtheorem{df}[thm]{定義}
\newtheorem{exam}[thm]{例}
\newtheorem{fact}[thm]{事実}
\newtheorem*{claim}{主張}
\newtheorem*{cau}{注意}
\renewcommand{\proofname}{{\bf 証明}}
%矢印たち。
%\toは\rightarrowと同じ。\getsは\leftarrowと同じ。同値は\iff
%上向き矢はuparrow逆はdownarrow
\newcommand{\To}{\Longrightarrow}
\newcommand{\Gets}{\Longleftarrow}
%黒板太字
\newcommand{\nn}{\mathbb{N}}
\newcommand{\zz}{\mathbb{Z}}
\newcommand{\qq}{\mathbb{Q}}
\newcommand{\rr}{\mathbb{R}}
\newcommand{\cc}{\mathbb{C}}
%無限大
\newcommand{\mgn}{\infty}
%ギリシャン
\newcommand{\dpst}{\displaystyle}
\newcommand{\ep}{\varepsilon}
\newcommand{\al}{\alpha}
\newcommand{\be}{\beta}
\newcommand{\de}{\delta}
\newcommand{\la}{\lambda}
\newcommand{\La}{\Lambda}
\newcommand{\ga}{\gamma}
\newcommand{\lil}{\la\in \La}
%商集合
\newcommand{\qt}[2]{#1/\! #2}
%enumerate環境
\renewcommand{\labelenumi}{{\rm (\arabic{enumi})}}
%以降alignじゃなくてenumerate を使え。
%なんか演算
\newcommand{\card}{\text{{\rm card}}}
\newcommand{\supp}{\text{{\rm supp}}}
%なんか演算２
\newcommand{\bd}{\text{{\rm Bdry}}}
\newcommand{\cl}{\text{{\rm CL}}}
\newcommand{\nai}{\text{{\rm Int}}}
\newcommand{\s}{\text{{\rm St}} }
%差集合は\setminusで統一する。等号の否定は\neq
%真部分集合は\subsetneqq　偏微分は\partial
%ディスプレイスタイル。\dfracでディスプレイ分数
%空集合は\Oを使う！！！！！！！！！！！！

\title{点可算被覆と可算コンパクト性}
\author{電波通信}
\date{\today}
			\begin{document}
\maketitle
この文書では位相空間に如何なる分離公理も課さない。
\begin{df}[星型集合]
集合$X$とその部分集合族$\mathcal{A}$、及び$x\in X$について
\[
\s(x,\mathcal{A})=\bigcup\{A\ |\ A\in \mathcal{A},\ x\in A\}
\]
を$\mathcal{A}$による$x$を中心とする星型集合という。
\end{df}

\begin{df}[完全集積点]
位相空間$X$の部分集合$A$について、$x\in X$が$A$の完全集積点であるとは、$x$の任意の開近傍$U$について
\[
\card(A)=\card(U\cap A)
\]
となること。
\end{df}

\begin{thm}\label{aa}
位相空間$X$について以下は同値である。
\begin{enumerate}
\item $X$は可算コンパクト
\item $X$の可算部分集合$A$に完全集積点が存在する。
\item $X$の無限部分集合$A$について点$a\in X$が存在して$a$の任意の近傍$N$について
\[
\card(A\cap N)\ge \aleph_0
\]
を充たす。
\end{enumerate}
\end{thm}
\begin{proof}
$(2)$と$(3)$が同値であることはよくわかる。$(1)$と$(2)$が同値であることを示そう。対偶を証明する。\par
[$\lnot (2)\To \lnot(1)$の証明]\par
$\lnot(2)$より、完全集積点を持たない可算無限集合$A$が存在する。完全集積点を持たないので各点$x\in X$について$x$の近傍$U_x$が存在して$\card(U_x\cap A)<\aleph_0$となる。さて$A$の有限部分集合$S$について
\[
O_S=\bigcup\{U_x\mid U_x\cap A=S\}
\]
と定義しよう。$A$の有限部分集合全体$\mathcal{F}$は可算集合であるから、$\{O_S\}_{S\in \mathcal{F}}$は可算な族であり、$X$の開被覆である。しかし構成の仕方からわかるように$\{O_S\}_{S\in \mathcal{F}}$の有限部分被覆$\mathcal{U}$の和集合$\bigcup\mathcal{U}$と$A$について$(\bigcup\mathcal{U})\cap A$は有限集合である。よって$X\neq \bigcup\mathcal{U}$。つまり$\{O_S\}_{S\in \mathcal{F}}$は有限部分被覆を持たない$X$の可算開被覆である。ゆえに$X$は可算コンパクトでない。\par
[$\lnot (2)\To \lnot(1)$終わり]\par
[$\lnot (1)\To \lnot (2)$の証明]\par
$X$は可算コンパクトでないから有限部分被覆を持たない可算開被覆$\mathcal{U}=\{U_i\}_{i\in \nn}$が存在する。そこで$\{a_i\}_{i\in \nn}$を
\[
a_n\not\in \bigcup_{i=0}^nU_i
\]
でかつ$n\neq m\To a_n\neq a_m$となるように選ぶ（これはちゃんと実行できる）。そして$A=\{a_n\}_{n\in \nn}$とすると$A$は可算無限集合である。これが完全集積点を持たないことを示そう。任意の$x\in X$について$x\in U_m$となる$\mathcal{U}$の元をとると、構成の仕方から$U_m$は高々$m+1$個の$A$の元としか交わらない。よって$x$は任意であったから$A$は完全集積点を持たない$X$の可算部分集合である。\par
[$\lnot (1)\To \lnot (2)$終わり]
\end{proof}

\begin{thm}
可算コンパクト空間$X$の点可算な開被覆$\mathcal{U}$は有限部分被覆を持つ。
\end{thm}
\begin{proof}
$\mathcal{U}$を点可算な$X$の開被覆としよう。
\[
\mathcal{M}=\{M\subset X\mid \forall a\in M\ \s(a,\mathcal{U})\cap M=\{a\}\}
\]
と定義する。このとき$\mathcal{M}$は非空でかつ包含関係についてZornian ordered set\footnote{Zornの補題の仮定を充たす順序集合のことをZornian ordered setとここでは呼ぶことにする。}である。よって極大元が存在する。その一つを$M$としよう。$M$は有限集合である。もしそうでないとしたならば、可算コンパクト性と定理\ref{aa}から$p\in X$が存在し$p$の任意の近傍と$M$との共通部分が無限集合になるが、$p\in U\in \mathcal{U}$となる$U$を取ると、$M$の取り方から$U$は$M$と高々一点でしか交わらないので矛盾する。\par
さて$M$の極大性から
\[
X=\bigcup_{a\in M}\s(a,\mathcal{U})
\]
となる。つまり$\mathcal{V}=\{V\in \mathcal{U}\mid V\cap M\neq\O\}$は$X$の被覆である。さらに点可算性から$\mathcal{V}$は高々可算であるので$X$の可算コンパクト性から$\mathcal{V}$の有限部分被覆が存在するが、$\mathcal{V}\subset \mathcal{U}$であるからそれは$\mathcal{U}$の有限部分被覆でもある。
\end{proof}

\begin{df}[メタリンデレフ]
位相空間$X$の任意の開被覆に点可算な開細分被覆が存在するとき、$X$はメタリンデレフであるという。
\end{df}

\begin{cor}
メタリンデレフかつ可算コンパクトな空間はコンパクトである。
\end{cor}


\begin{cor}[Arens-Dugundji]
メタコンパクトかつ可算コンパクトな空間はコンパクトである。
\end{cor}


\begin{cor}
$\omega_1$はメタリンデレフではない。より詳しく、$\omega_1$の有界開集合からなる開被覆は点可算な細分を持たない。
\end{cor}






































	
\begin{thebibliography}{9}
\bibitem{1}L.A.Steen and J.A.Seebch,\ {\it Counterexamples in Topology},\ Dover Publications Inc New York,1995
\end{thebibliography}




			\end{document}



\begin{comment}
[集合のやつは$\mid$を使う。]
[真なる包含関係]
\subsetneqq
[可換図式]
\[\xymatrix{
&   & (2) \ar@{=>}[rd]&  &  &  & (5) \ar@{=>}[rd]&\\ 
& (1)\ar@{=>}[ru] \ar@{=>}[rd]&  &(4)\ar@{=>}[r]&(7)& (1)\ar@{=>}[ru] \ar@{=>}[rd]& & (7)&\\ 
&   & (3) \ar@{=>}[ur]&  &  &  &(6) \ar@{=>}[ur]&
}\]

[\bigin \endを使わない方法]

{\prop[aa] aaa}
で
\begin{prop}[aa]
aaa
\end{prop}
と同じ\df \thm \proof でも同様

[直和]
$\amalg$

[同値関係でよく使う波線]
\sim

[引数付きの新命令]
\newcommand{\prn}[1]{\|#1\|_p}

[番号のリセット]

\setcounter{equation}{0}


[字体の変更]

{\rm hogehoge}と使う。

\rm ローマン
\it イタリック
\sl 斜体

[supp]

\newcommand{\supp}{\text{{\rm supp}}}

[中央揃えの改行]

	\begin{eqnarray*}
		hoge1&=&hogehoge
		hoge2&=&hogehofe

	\end{eqnarray*}

&&の部分で揃えられる。


[\tag]
\begin{equation}
f(x) = ax^2 + bx + c \tag{1.2.3}
\end{equation}



[場合分け]

	\begin{cases}
		hoge1 & \text{状況１}\\
		hoge2 & \text{状況２}\\
		・
		・
		・
	\end{cases}



\end{comment}





